\chapter{Troubleshooting and Frequently Asked Questions}\label{chap:12}
\label{ch:trouble}

This chapter collects the messages that REGEN can show, explains why they appear and what to do. Messages are
quoted as the program shows them; parts in angle brackets (such as $\langle$name$\rangle$) are filled in with the
actual factor, response or number. The last sections answer common questions about reproducibility, speed and
installation.

\section{Installation, starting and files}
\label{sec:ch12-files}

{\small\begin{longtable}{@{}p{0.38\linewidth}p{0.58\linewidth}@{}}
\caption{Installation and file messages.}\label{tab:ch12-files}\\
\toprule
Message or problem & Cause and what to do \\
\midrule
\endfirsthead
\toprule
Message or problem & Cause and what to do \\
\midrule
\endhead
\bottomrule
\endfoot
``Windows protected your PC'' when the installer starts &
Windows SmartScreen warns about a downloaded program that is not code-signed. Check that the file is the REGEN
installer from the GitHub Releases page, click \ui{More info} and then \ui{Run anyway}
(Section~\ref{sec:ch02-install}). \\
\addlinespace
The installer asks for administrator rights &
\ui{Install for all users} was chosen. Choose \ui{Install for me only}, which needs no administrator rights. \\
\addlinespace
Double-clicking a \texttt{.doe} file does not open REGEN &
The file association was not selected during installation. Run the installer again and tick \ui{Open .doe files
with REGEN}, or open the file with \menupath{File $\rightarrow$ Open\dots}. \\
\addlinespace
``Open Failed: The file could not be opened: $\langle$error$\rangle$'' &
The file is not a REGEN/DOE Studio project, it was damaged (for example cut off while copying, or edited by hand
and no longer valid JSON), or it is locked by another program. Open a backup copy; the error text names the
problem (for example the line of a JSON error). \\
\addlinespace
``Save Failed: $\langle$error$\rangle$'' &
The folder is read-only, the disk is full, or the file is open in another program (for example a cloud-sync
client). Save to another folder with \ui{Save As\dots}. \\
\addlinespace
``Export Failed'' / ``PDF Export Failed'' &
The target file is open in another program (an Excel workbook or PDF viewer keeps the file locked) or the folder is
not writable. Close the file or choose another name. \\
\addlinespace
``Export data: The file could not be written: $\langle$error$\rangle$'' &
Same as above, for \ui{Export data\dots}, \ui{Export table\dots} and \ui{Export all tables\dots}. \\
\addlinespace
``Export data: There are no data to export.'' &
The plot or panel is empty (for example a message is shown instead of a plot). \\
\addlinespace
A project opens without the last NSGA result or search table &
Results larger than 5 million characters (for example a very large Pareto history) are not stored in the file, and
projects saved by older versions contain no results. Run the analysis again. \\
\addlinespace
``User Manual: The user manual (REGEN\_User\_Manual.pdf) was not found next to the program\dots'' &
The PDF manual is missing from the installation folder. Reinstall REGEN or download the manual from the Releases
page. \\
\end{longtable}}

\section{New design and import}
\label{sec:ch12-design}

{\small\begin{longtable}{@{}p{0.38\linewidth}p{0.58\linewidth}@{}}
\caption{Messages when creating a design or importing data.}\label{tab:ch12-design}\\
\toprule
Message & Cause and what to do \\
\midrule
\endfirsthead
\toprule
Message & Cause and what to do \\
\midrule
\endhead
\bottomrule
\endfoot
``Factor $\langle$name$\rangle$: Low must be less than High.'' / ``The limits of $\langle$name$\rangle$ are not
numbers.'' & Correct the factor table. Decimal comma and point are both accepted. \\
\addlinespace
``The sum of the component lower limits must be less than the total.'' &
In a mixture the lower limits must leave room for the components to vary: $\sum L_i < T$. Lower some limits or
raise the mixture total. \\
\addlinespace
``Categorical factor $\langle$name$\rangle$ needs at least 2 distinct levels.'' &
Separate the level names with semicolons (or commas), for example \texttt{Ethanol; Methanol}. \\
\addlinespace
``Constraint: Line $\langle i\rangle$: A constraint must use <= or >=.'' and similar &
A linear constraint could not be read. Write one constraint per line with factor letters, numbers and \texttt{<=}
or \texttt{>=}, for example \texttt{A + B <= 160} or \texttt{2A - C >= 0}. Constraints may not contain categoric
factors, and only letters of existing factors (``Factor 'X' does not exist (use letters A–C).''). \\
\addlinespace
``The design region is empty: the bounds and constraints contradict each other.'' &
No point satisfies all limits and constraints at the same time. Relax a constraint or widen a factor range. \\
\addlinespace
``There are not enough candidate points to estimate this model. Relax the constraints or choose a simpler model.'' &
Optimal designs choose runs from a candidate set; in a small constrained region the set cannot support all model
coefficients. Use a lower-order planned model or relax the constraints. \\
\addlinespace
``Box-Behnken requires at least 3 factors.'', ``DSD is available for 3–20 factors.'', ``PB with $N$ runs allows at
most $N-1$ factors.'', ``$\langle$array$\rangle$ allows at most $m$ factors.'' &
The chosen design type does not exist for this number of factors. Choose another type or size (the Design Guide
lists the valid choices). \\
\addlinespace
``Split-plot: Enter the letters of the hard-to-change factors (e.g. A).'' &
The split-plot option is ticked but no hard-to-change factor is given. \\
\addlinespace
\ui{Augment Design\dots} is greyed out; tooltip ``Augmenting split-plot designs is not supported: the new runs would
need their own whole plots. Create a new split-plot design instead.'' &
Split-plot designs cannot be augmented. Create a larger split-plot design, or add runs suggested by Data Readiness
(each new run becomes its own whole plot). \\
\addlinespace
``Import: The clipboard is empty. Copy a data table first.'' &
Nothing was copied, or the copied content is not text. Copy the cells again in Excel. \\
\addlinespace
``Mark at least one Factor column and one Response column.'' &
Set the column roles by right-clicking the column headers (Section~\ref{sec:ch02-import}). \\
\addlinespace
``Factor $\langle$name$\rangle$ is constant / empty - mark it as Ignore.'' &
A factor must take at least two different values. Constant columns carry no information for a model. \\
\addlinespace
``Column $\langle$name$\rangle$, row $i$ is not a number: '$\langle$value$\rangle$'.'' &
A cell of a numeric column contains text (such as ``n/a'' or a unit). Delete the row, clear the cell before
copying, or mark the column as categoric. Cells shown in red are the ones that will fail. \\
\addlinespace
``$\langle$name$\rangle$ has $L$ levels - should this column be numeric or ignored?'' &
A text column with more than 20 different values was marked as a categoric factor. It is probably an identifier
(mark it \ui{Ignore}) or a numeric column with some text cells. \\
\addlinespace
``Mixture components cannot be combined with process factors when importing data.\dots'' &
An import is either a process or a mixture data set. Mark all factor columns as components, or none. \\
\addlinespace
``In $n$ rows the components do not sum to the total $T$. Continue (those rows are still analyzed)?'' &
Check \ui{Mixture total}; it is set to the median row sum. Small rounding differences are harmless; large ones
indicate a wrong column or a wrong total. \\
\addlinespace
``Sum of component minimums ≥ total - check the mixture total.'' &
The lower limits taken from the data leave no room for variation. Check the total and the component columns. \\
\addlinespace
Numbers such as 1,234 are read as 1.234 &
A comma alone is read as a decimal comma. Remove thousands separators before copying. \\
\end{longtable}}

\section{Data, analysis and graphs}
\label{sec:ch12-analysis}

{\small\begin{longtable}{@{}p{0.38\linewidth}p{0.58\linewidth}@{}}
\caption{Messages on the data, analysis and graph pages.}\label{tab:ch12-analysis}\\
\toprule
Message & Cause and what to do \\
\midrule
\endfirsthead
\toprule
Message & Cause and what to do \\
\midrule
\endhead
\bottomrule
\endfoot
``The model cannot be computed yet'' with ``Fill in at least 3 complete rows (factors + response) to start the
analysis.'' & The response has fewer than 3 runs with complete factor settings and a response value (excluded runs
do not count). Enter more data. \\
\addlinespace
``Residual degrees of freedom = 0 (saturated model). Statistical tests cannot be computed - remove model terms or
add runs.'' &
The model has as many coefficients as there are runs, so no error estimate exists. Choose a lower-order model,
remove terms, or add runs (Data Readiness suggests D-optimal runs). Design Evaluation shows the degrees of freedom
before any run is made. \\
\addlinespace
``Aliased terms (removed automatically): $\langle$terms$\rangle$'' &
These terms cannot be estimated separately from terms earlier in the model (for example a cubic term in a
Box-Behnken design, or interactions in a resolution III fraction). They are dropped from the model; the Fit Summary
marks such orders ``Aliased''. Use a design of higher resolution, add runs (augmentation, foldover) or use a
lower-order model. \\
\addlinespace
``No model is significant compared with the mean.'' &
None of the model orders explains the response better than its mean (sequential $p \ge 0.05$). The factors may
have no effect in the studied range, the noise may be large, or the ranges may be too narrow. \\
\addlinespace
``The square root transform requires all response values ≥ 0.'' / ``The $\langle$transform$\rangle$ transform
requires all response values > 0.'' &
Set \ui{Constant k (y + k)} so that $y+k$ is positive, or choose another transformation. \\
\addlinespace
``The Box-Cox plot needs all values of y + k > 0 (and at least 3 data points).'' / ``Box-Cox is not available for
this response.'' &
Box-Cox needs positive values and is only available for ordinary least squares (not for logistic, Poisson or
split-plot models). \\
\addlinespace
``No effect exceeds the ME limit.'' &
Lenth's automatic selection found no effect larger than the margin of error. Inspect the half-normal plot; select
terms by clicking if there is a clear break. \\
\addlinespace
``The GLM fit has warnings - see the GLM (Logistic/Poisson) tab.'' &
The IRLS fit did not converge, or fitted probabilities are numerically 0 or 1 (complete or quasi-complete
separation; for Poisson: fitted means 0). Standard errors, tests and odds ratios are then not reliable. Remove terms
or add runs where the response changes. \\
\addlinespace
``A binomial response must be a number of successes between 0 and $m$ (number of trials).'' / ``A Poisson
response (counts) cannot be negative.'' &
Check the response type and \ui{Trials per run} in \ui{Factor \& Response Info\dots}, and the data. \\
\addlinespace
``Whole-plot df = 0: hard-to-change factor effects cannot be tested properly. Add more whole plots (replicates).'' &
In a split-plot design the hard-to-change factors are tested against the whole plots. With too few whole plots no
test is possible. \\
\addlinespace
``This slice lies entirely outside the design region. Change the values of the other factors/components.'' &
With the values chosen in \ui{Other Factor Values} (or \ui{Other Components}), no point of the plotted slice
satisfies the constraints or component bounds. Move the other factors toward the center of the region. \\
\addlinespace
``This graph needs more numeric factors.'' &
The graph type needs two (contour, surface, overlay) or three (cube) numeric factors. \\
\addlinespace
``No criteria yet. Choose response goals and limits on the Desirability page.'' &
Overlay and desirability contours use the goals of the Desirability page. Set at least one goal there. \\
\addlinespace
``Enter the factor standard deviations first (Data page $\rightarrow$ Factor \& Response Info\dots).'' &
Propagation of error needs the standard deviation of at least one factor. \\
\addlinespace
``No complete runs to evaluate yet.'' (Design Evaluation) &
All runs have missing factor settings or are excluded. \\
\addlinespace
``Simulation: Invalid equation: \dots'' / ``Unknown name 'x'.\dots'' &
Use the factor letters (A, B, \dots), \texttt{+ - * / **} (or \texttt{\^{}}), and the listed functions; use a point
as decimal separator inside an equation (commas separate the arguments of \texttt{min} and \texttt{max}). \\
\end{longtable}}

\section{Artificial neural networks}
\label{sec:ch12-ann}

{\small\begin{longtable}{@{}p{0.38\linewidth}p{0.58\linewidth}@{}}
\caption{Messages on the ANN page.}\label{tab:ch12-ann}\\
\toprule
Message & Cause and what to do \\
\midrule
\endfirsthead
\toprule
Message & Cause and what to do \\
\midrule
\endhead
\bottomrule
\endfoot
``ANN here is for continuous (normal) responses. For binomial/Poisson responses use logistic/Poisson regression on
the RSM Analysis page.'' & ANN models are trained only for normal responses. \\
\addlinespace
``Too little data to train an ANN (at least 4 complete runs required).'' / ``Too little training data. Increase the
training data percentage.'' & Enter more data or change \ui{Training data:} so that the training set is not empty. \\
\addlinespace
``Training + validation + test percentages must add up to 100\%.'' & Adjust the three percentages. \\
\addlinespace
``Early stopping requires validation data (> 0\%).'' &
Set \ui{Validation data:} above 0, or untick early stopping. \\
\addlinespace
``Enter the neurons per hidden layer as numbers separated by commas. Example: 20, 10, 5'' &
One number per hidden layer; each layer needs at least 1 neuron. \\
\addlinespace
``Constant response - the ANN cannot be trained.'' &
All training values are equal; there is nothing to learn. \\
\addlinespace
``Only $N$ complete runs. Architecture search needs at least 8 runs.'' &
Cross-validation needs enough runs to form folds. Use a fixed architecture or RSM. \\
\addlinespace
``The learning curve needs at least 12 complete runs (found $N$).'' &
With fewer runs no learning curve can be estimated. \\
\addlinespace
Learning curve: ``Too little data ($N$ runs) for a reliable learning curve.'' &
The curve was computed but $N<25$, so its extrapolation is not trusted. An ANN generally needs at least 20 to 30
runs; use RSM. \\
\addlinespace
Learning curve: ``(run estimate unreliable)'' &
The fitted power law reached the edge of the exponent range (0.2 to 2.5), its error limit is 0, or the
interquartile range of the run estimate is very wide (upper/lower above 3). Use the curve only to decide whether
more runs help, not how many. \\
\addlinespace
Status ``Overfitting'' or ``Underfitting'' &
Overfitting: the error on new data is far above the training error; use fewer neurons, Bayesian regularization,
or more data. Underfitting: even the training error is high; use more neurons or check the data. After an
architecture search the headline status comes from cross-validation; the single train/test split of the final
model is shown separately and is less reliable with few runs. \\
\addlinespace
Negative ANN or RSM test $R^2$ &
With few runs a test set of 2--3 runs gives a meaningless $R^2$. Use the cross-validated $Q^2$ of the
architecture search. A negative RSM cross-validation $R^2$ for Box-Behnken or CCD data is explained in the report:
leaving runs out destroys the design structure; compare with the RSM predicted $R^2$. \\
\addlinespace
``The search was stopped before finishing; the results above are from the architectures evaluated so far.'' &
\ui{Stop} was clicked. The table and the selection use only the finished candidates. \\
\addlinespace
``Architecture $\langle$layers$\rangle$ is too large for the computer's memory; adding neurons was stopped.'' /
``A parallel process stopped unexpectedly (possibly not enough memory).'' &
Reduce \ui{CPU cores (parallel):}, set a maximum number of neurons, or close other programs. \\
\addlinespace
``Large data ($N$ runs): the search uses 4000 random runs for speed; the final model is trained on all data.'' &
Information only. \\
\end{longtable}}

\section{Optimization, sensitivity and prediction}
\label{sec:ch12-opt}

{\small\begin{longtable}{@{}p{0.38\linewidth}p{0.58\linewidth}@{}}
\caption{Messages on the optimization, sensitivity and prediction pages.}\label{tab:ch12-opt}\\
\toprule
Message & Cause and what to do \\
\midrule
\endfirsthead
\toprule
Message & Cause and what to do \\
\midrule
\endhead
\bottomrule
\endfoot
``Choose a goal (Maximize/Minimize/Target/In range) for at least one response whose model is available.'' &
All response goals are \ui{None}, or the responses with a goal have no model yet. Set a goal in the
\ui{Goal} column. \\
\addlinespace
``Skipped responses without a model: $\langle$names$\rangle$'' &
These responses have a goal but no model (not enough data, or ANN chosen but not trained). The optimization runs
without them. \\
\addlinespace
``Limits for response $\langle$name$\rangle$: lower must be < upper.'' / ``The target for $\langle$name$\rangle$
must be between the limits.'' & Correct the limits in the criteria table. \ui{Reset Limits to Data} restores the
data minimum and maximum. \\
\addlinespace
``No point satisfies all limits and constraints.'' / ``The component limits do not allow the mixture to sum to the
total.'' &
The factor limits on the Desirability page, the component bounds and the constraints leave no feasible point. Widen
the factor limits. \\
\addlinespace
Desirability $D = 0$ for all solutions &
At least one response cannot reach its lower limit (maximize) or upper limit (minimize) anywhere in the region.
Relax that limit; the bar chart of a solution shows which individual $d_i$ is 0. \\
\addlinespace
``Select at least 2 objectives (maximize/minimize/target).'' / ``NSGA-II requires at least 2 objectives\dots'' &
NSGA-II/III builds a Pareto front between two or more objectives. For one objective use the Desirability page. A
row with goal \ui{Ignore} but with limits counts only as a constraint. \\
\addlinespace
``Enter a target value for $\langle$name$\rangle$.'' & Goal \ui{Target} needs a value in the \ui{Target} column. \\
\addlinespace
``The model for $\langle$name$\rangle$ is not available yet.'' & Fit the RSM model or train the ANN for that
response first. \\
\addlinespace
``No feasible solution. Relax the response limits or constraints.'' &
No solution of the final population satisfies the response limits. Widen the limits or remove some. \\
\addlinespace
``There is no ANN model for the objective responses yet. Train an ANN first on the ANN page (Optimize All Responses
button).'' & \ui{Compare ANN vs RSM} needs an ANN for the objectives. \\
\addlinespace
``Run the optimization first.'' & \ui{Check Population Size} and \ui{Send Selected Solution to Prediction} need a
result. \\
\addlinespace
``Automatic stop was not used: with NSGA-II and 4 or more objectives\dots'' &
For NSGA-II with four or more objectives the automatic stop is switched off and all generations are run. Use
NSGA-III (or \ui{Automatic}) for a reliable automatic stop. \\
\addlinespace
``Not enough generations yet\dots'' / ``Population: not sufficient\dots'' &
Check C was not passed. Increase \ui{Max generations:} or \ui{Population:} and run again
(Table~\ref{tab:ch01-checks}). \\
\addlinespace
``The models differ considerably\dots'' (red) &
Check D: the RSM and ANN fronts disagree. Make confirmation runs at both TOPSIS points and trust the model with the
better predictive accuracy. \\
\addlinespace
``Sobol indices need at least one process or categoric factor. Mixture components sum to a constant and are not
independent inputs.'' &
Sobol indices, standardized coefficients and the tornado chart vary independent factors only. In a pure mixture use
the parametric study, the trace plot or the Pareto trade-off instead. \\
\addlinespace
``The linear constraints of the design are not applied here\dots'' &
Sobol indices assume independent factors, so they describe the whole factor box, including points outside the
constraints. \\
\addlinespace
``Too few points of the sample satisfy the design constraints.'' &
The constrained region is too small a part of the factor box for the Latin hypercube sample. Increase
\ui{Sample size:}. \\
\addlinespace
``Reference point: The reference point lies outside the design region (component bounds or constraints).\dots'' &
Change the reference values or choose another start for the sensitivity study. \\
\addlinespace
``Run NSGA-II / NSGA-III on the multi-objective optimization page first.'' & The Pareto trade-off analysis needs an
NSGA result. \\
\addlinespace
``The point is outside the design region - the prediction is an extrapolation.'' &
A factor is set beyond the range of the design and of the data. Predictions there are not supported by data. \\
\addlinespace
``This point violates the component bounds / constraints.'' / ``The component sum does not equal the mixture
total.'' & Correct the prediction point. \\
\addlinespace
Prediction shows ``-'' for the standard error, CI and PI &
The response uses an ANN, which gives only a point prediction. Confirmation is then judged by the error as a
percentage of the response data range (at most 5\,\%). \\
\end{longtable}}

\section{Frequently asked questions}
\label{sec:ch12-faq}

\paragraph{Why do I get a different run order every time I create the same design?}
The run order (and the choice of points in optimal, Latin hypercube and hybrid designs) is random. The
\ui{Random seed:} in the New Design dialog controls it: the same seed gives the same design. A new random seed is
proposed every time the dialog opens; type a fixed number to reproduce a design. The seed is saved with the project
and shown on the Summary page and in the PDF report.

\paragraph{Are the results reproducible?}
Yes. Every random step has a seed: the design seed, \ui{Random seed:} on the ANN page (data split, initial weights,
cross-validation folds), \ui{Seed:} on the NSGA-II / III page, and fixed internal seeds for the desirability
search, the region samples of Data Readiness and Design Evaluation, and the sensitivity samples. With the same data,
settings and seeds the same numbers are obtained. ANN results do not depend on the number of CPU cores.

\paragraph{The ANN gives a different result each time I click Train ANN.}
This happens only when settings change between clicks (data, split percentages, seed, algorithm). Keep the seed
fixed; to see how much the result depends on the initialization, change only the seed.

\paragraph{The calculation is slow. What can I do?}
For ANN work increase \ui{CPU cores (parallel):} (the default already uses all cores but one), reduce
\ui{Restarts (initializations):} or \ui{Maximum epochs:}, or limit the number of neurons and layers of the
architecture search. With more than 5000 runs the number of restarts of the final training is limited
automatically, and with more than 4000 runs the architecture search uses a random subset of 4000 runs. For
NSGA-II/III keep \ui{Stop automatically on convergence} on. For RSM, data sets with more than 400 runs are handled
with a faster algorithm automatically; data sets with 15\,000 runs have been tested.

\paragraph{The computer becomes slow while REGEN trains networks.}
Reduce \ui{CPU cores (parallel):}. The worker processes close by themselves 90 seconds after the last calculation.

\paragraph{Can I use REGEN without administrator rights?}
Yes: choose \ui{Install for me only} during installation.

\paragraph{Where is the list of recent projects stored?}
In the Windows user settings (registry) under the old program name ``DOE Studio'', so the list of earlier versions
is kept. Deleted or moved files are not shown.

\paragraph{Can I change the factor ranges of an existing design?}
The factor settings of a designed experiment are fixed once the design is created (the data sheet shows them in
grey). For historical data (created with Import Data or the Historical Data design types) the factor cells can be
edited on the Data page.

\paragraph{How do I remove an outlier without deleting it?}
Select the run on the Data page and click \ui{Exclude/Restore Runs}, or select the point on the diagnostics plots
and click \ui{Exclude Selected Runs}. Excluded runs are left out of all models, Design Evaluation and Data
Normalization, but stay in the data sheet and in the exports (marked ``Excluded'') and can be restored at any time.
